\documentclass[12pt]{article}

\usepackage[a4paper, total={6in, 8in}]{geometry}
\usepackage{textcomp}

\begin{document}
\setlength{\parindent}{0pt}

\def \t {\textrightarrow}
\def \v {\vspace{1em}}
\def \bi {\begin{itemize}}
\def \ei {\end{itemize}}
\def \s[#1] {\section*{#1}}
\def \ss[#1] {\subsection*{#1}}
\def \sss[#1] {\subsubsection*{#1}}

\s[D'annunzio]
Forte connessione tra vita e arte \t estetismo \t la poesia e il focus su Alcyone \\
Forza di vita in "Canta la gioia" \t c'è pero una fase di tornare alla casa, alla famiglia = ai valori semplici \t non esiste solo il d'annunzio di Alcyone, ma il d'annunzio di "Poema paradisiaco" e "Notturno" \\
Eiste anche un d'annunzio piu intimo, piu profondo \t quindi + assertivo: questo è il d annunzio che descrive il sentore decadente, il d'annunzio europeo, che caratterizza una frangia del decadentismo: i crepuscolari \\
Questa è la voce di d'annunzio che resta \t il ruolo del poeta \\
Luciano Angeschi \t il poeta perde l'aureola \t se qui il poeta è il vate, poi si domandera a cosa serve la sua voce, se la voce dell'anima rompe le barriere per farsi sentire ?? \\
"Notturno" \t alla luce di canta la gioia, qua l'autore preferisce la notte \t un ripiegamento intimo e profondo \\
È presente una ricognizione di tutti i rapporti umani \t quali sono le figure che costruiscono l'essere umano? da subito la famiglia, gli amici \\
Nel notturno, d'annunzio non vede con gli occhi ma con il cuore \\
Il "Poema paradisiaco" \t "Consolazione" è come un romanzo in versi, una prosa versificata \t come "L'addio ai monti di manzoni" \t in "consolazione" dice che vuole tornare dalla madre, e questo mette a luce un meccanismo freudiano: esplosione di eros, tanatos \\
Camillo Sbarbaro: "Taci anima stanca di godere" \t disfattismo interiore che accompagna una fase di intenso vitalismo, che depaupera la persona di energie e volonta \\ 
L'unico modo per tornare a casa è quello di abbracciare gli affetti semplici \\
Volto + decadente e + europeo \\
Paradisiaco riporta all'eden: 3 riferimenti
\bi
    \item all hortus conclusus \t è l'indicazione di un giardino chiuso
    \item hortus dedicato alle leggerezze
    \item hortus + significativo è l'hortus animae
\ei

\v 

Ricordiamo "Mille e mille" \t "Il libro segreto" di d annunzio, in cui c'è una confessione privata dei sui gusti sentimentali e delle sue relazioni \\
"Discorso di Quarto" \\
Lui vive perennemente indebitato \\
Capponcina, colli della toscana \t è un luogo dove si reca e vive vita intensa con Eleonra Aduse \\
Papini e Prezzolini, "Il bagno di sangue" \t spirito bellicoso e voyerismo, vitalismo d annunziano \t Bava Beccaris \\
Nel ?? il Tovez, sulla gazzetta letteraria, parla della tendenza di plagio e copiatura, poco di fine ripresa dei versi di altri nella poesia di d'annunzio \\
Luciano Angeschi parla della perdita dell'aureola, e parla di angoscia nascosta dietro il voyerismo dell'autore \t lui riempie il vuoto angoscioso con questi oggetti \\
Sono modi per riempire un vuoto \t ed è il contrario del minimalismo di altri autori, ovvero la presenza esistenzialista nella poesia \t qua invece voyerismo marcato, e gli spazi sono riempiti dagli oggetti, che colmano una solitudine \\
Questo è tipico di un uomo decadentista e un mentalista estetizzante \\
In poesia si traduce in stile pleonastico, utilizzo di termini aulici, sovrabbondanza di nomi, e utilizzo della parola come in "Art poetique" di Verlain, che parla dell'arte poetica \\
Si puo dire che d'annunzio si presente come il poeta dell'eccesso, e l'eccesso si deve leggere nell'esaltazione pagana della vita \t che troviamo in "Canta la gioia" \\
Poesia dell'eccesso, della sovrabbondanza

\ss[Confronto con pascoli]
Non vive a contatto con il mondo, il suo modo di essere fuori dagli schemi è di entrare nel fanciullino \t pascoli si rivolge all interno, mentre d'annunzio all'esterno, nella sovrabbondanza degli oggetti \\
Cosi entrambi si differenziano dalle masse \\
Anche il vittoriale non ha praticamente spazi vuoti \t tiene un letto che non utilizza ???, dove lascia camice da notte = spazialita e liberta dell individuo \\
In d annunzio la spazialita è riempita, è l'esempio dell suo animo \t per riscattarsi dalla prigionia del bello, si circonda di oggetti che lo avvolgano \\
È il vuto dell'artista decadente \t qua c'è il significato del limite, che è l'ortus conclusus \t e poi si confrontera con l hortus montaliano \\
Per superare d'annunzio occorre attraversarlo

\ss[Ascetismo ed estetismo]
Si trovano coniugati nell'autore \t parabola di vita complessa \t ha fasi di depressione e sofferenza profonda \\
È esteta per riempire l'horror vacui \t non sceglie ma affoga negli oggetti e nella bellezza \t allora non prende le distanze dall'inettituinde: non si riesce a riprendere dalla sovrabbondanza di oggetti \\
Vita senza sentimento non è garanzia di felicita \\
Ha solo valenza riempitiva \\
La sua personalita estetizzante segue la strada del riempire con cio che trova \t Aduse diventa riempimento per un periodo \t e da spunto per opere teatrali ??? \\
Lui è il prototipo del fallimento \t il piacere è come il canto V dell'infermo: la celebrazione dell'amore passionale porta al fallimento \t Andreas Perelli rappresenta il voyerismo e il vitalismo e allo stesso tempo ne rappresenta il fallimento \t lui muore solo, non c'è la Muti e non c'è Maria \\
Andreas perelli capisce che nelle due donne ci sono due sfaccettature dell amore (maria avvolgente e accudente, mentre elena estremizzazione della sessualita) \t ma lui perde sia l una sia la tra \\
È una parabola che si torce su se stessa \t Paolo e Francesca stimatizzano la quintessenza dello stil novo, pero l'amore li ha portati alla morte \t qui le passioni sfrenate, la mania estetizzante e rifugiarsi nella bellezza e sensualita porta alla rovina \\
La parabola di perelli trova nel vertice il declino \t il suo essere artefice del vivere inmitabile, essere il nerbo del piacere: lo porta a sperimentare che l'assenza di un senso profondo è cio che lo candida alla solitudine \\
Le ragioni dell'innamoramento non ci sono \\
Lui sceglie la muti sulla base delle sfaccettature forti \t maria è accudente, emotiva, avvolgente, sensibile (che pero non toglie la sensualita, instaura solo un altro canale) \t pero non gli basta \\
Il vuoto che non puo colmare è in realta dentro di se \t ha un vuoto interiore \\
Erik Fromm \t è diverso dire "ti amo perche ho bisogno di te" e il contrario
\end{document}
